\chapter{结论与展望}


\section{研究结论}

总结上述分析结果，我们可以得到如下观点：

\begin{enumerate}[(1)]
  \item 总体而言，谈及教学设计技能，数学师范生掌握较好，但仍有进步空间.课堂教学技能也是如此，具体表现为语言技能和提问技能有待加强。 数学说课技能、教学评价技能和教学研究技能掌握情况较差,例如师范生难以从多个角度和不同的维度进行评课，查找、整理、分析文献能力较差，对学科发展动态和科研成果了解不足等。
  \item 学校对数学专业师范生的教学设计和课堂教学技能的培养较为重视，不论是教学设计次数和讲课次数，还是各项活动都对这两项教学技能的培养作用最大；忽视了数学说课、教学评价和教学研究技能的培养，导致师范生在技能发展上的不足。
  \item 师范生的数学教学技能的各个维度在年级上不存在明显差异；性别不同的师范生的数学教学技能水平存在明显差异，且女生的数学教学技能较男生好；普通话等级对数学师范生存在明显影响，普通话等级为一乙的师范生的数学教学技能掌握情况比普通话等级为二甲和二乙的师范生要好；计算机等级对师范生的数学教学技能影响不明显。
  \item 数学教学技能各个维度的之间存在明显的正相关，且课堂教学技能与教学评价技能相关性最大，数学说课技能与教学研究技能相关性最小；课堂教学技能的各个维度之间也存在明显的正相关，且结束技能和导入技能的相关性最高，讲解技能与媒体使用技能的相关性最低。
\end{enumerate}



\section{相关建议}

研究结果表明，高校数学师范生的教学设计和课堂教学技能的培养更加重视，在数学说课、教学评价和教学研究技能的培养方面存在不足，根据数学师范生的反馈和学院的培养计划，可以提出以下建议：

\begin{enumerate}[(1)]
  \item 提高数学说课技能的途径：教学比赛可以采取多种形式，讲课和说课相结合；微格教学除了训练学生的讲课和教学设计能力，也应当训练学生的说课能力；在理论学习的过程中，进行说课的指导和教学，给学生锻炼的机会。
  \item 提高教学评价技能的途径：在微格教学中，鼓励师范生相互评价，提高评价能力；在案例教学中，着重培养学生的教学评价能力，引导学生多角度、多方法进行评价；在教育见习和实习中，要求学生多参与实习学校的公开课听课和评课活动，在实践中提高能力。
  \item 提高教学研究技能的途径：除了现有的教学比赛，可以增加教学研究大赛，通过比赛培养学生的研究能力；学校可以开设相关选修课程，教授教学研究基本方法和步骤，并开展实践，在实践中培养学生的教学研究能力。
\end{enumerate}


\section{研究的不足与展望}

本研究存在许多不足与欠缺之处。首先是被试取样的范围比较局限，被试仅选取自华中师范大学数学与统计学学院，并且被试的男女比例不均衡，男生人数过少，无法显著代表数学师范生群体。今后，可以选取更多不同高校的数学师范生作为被试，并且尽量男女人数均衡，以进一步完善现有结论。其次，本次调查要求数学专业师范生对自身的教学技能掌握状况进行评判和评估，得出的结论为师范生自我报告的数学教学技能水平，研究结果存在一定主观因素影响。因此，在未来的研究中，可以从不同的视角对数学师范生的教学技能进行评分，例如学生、实习学校指导老师以及数学教育方面的专家等，使得研究更加客观、全面，以弥补现有研究的不足。